% META: {"id": "1SPE-TRIGO-EV-A-corrige", "chapitre": "1SPE-TRIGONOMETRIE", "type_objet": "corrige_evaluation", "evaluation_ref": "1SPE-TRIGO-EV-A", "version": "A", "capacites_codes": ["C1","C2"], "sources_inspiration": [], "parametres_sympy": {}, "fichier_tex": "chapitres/1SPE-TRIGONOMETRIE/evaluations/1SPE-TRIGO-EV-A-corrige.tex", "status": "needs_review", "approval_state": "STALE_AFTER_SEMANTIC_EDIT"}
% BEGIN-VERIFY
% from sympy import *
% x = symbols('x', real=True)
% assert 225 * pi / 180 == 5*pi/4
% assert 7*pi/6 * 180/pi == 210
% assert Rational(11,4) - 2 == Rational(3,4)
% assert cos(-2*pi/3) == Rational(-1,2)
% assert sin(-2*pi/3) == -sqrt(3)/2
% assert cos(5*pi/6) == -sqrt(3)/2
% assert sin(5*pi/6) == Rational(1,2)
% assert sin(7*pi/4) == -sqrt(2)/2
% assert Rational(3,5)**2 + Rational(4,5)**2 == 1
% assert 12 * 5*pi/6 == 10*pi
% assert Rational(9, 6) == Rational(3, 2)
% assert (7*pi/6) / (2*pi) == Rational(7, 12)
% assert cos(-pi/3) == Rational(1,2) and sin(-pi/3) == -sqrt(3)/2
% assert cos(3*pi/4) == -sqrt(2)/2 and sin(3*pi/4) == sqrt(2)/2
% assert 27 + 25*sin(pi/3) == 27 + 25*sqrt(3)/2
% assert abs(float(27 + 25*sqrt(3)/2) - 48.7) < 0.05
% END-VERIFY

\begin{corrige}{1SPE-TRIGO-EV-A}

%===============================================================
\section*{Exercice 1 \hfill (4 points) — C1}
%===============================================================

\baremeIndicatif{Q1 : 1 pt ; Q2 : 1 pt ; Q3 : 1 pt ; Q4 : 1 pt}

\textbf{1.} $225° = 225 \times \dfrac{\pi}{180} = \dfrac{5\pi}{4}$.

\textbf{2.} $\dfrac{7\pi}{6} = \dfrac{7\pi}{6} \times \dfrac{180}{\pi} = 210°$.

\textbf{3.} $\dfrac{11\pi}{4} - 2\pi = \dfrac{11\pi}{4} - \dfrac{8\pi}{4} = \dfrac{3\pi}{4} \in {]-\pi;\pi]}$. Mesure principale : $\dfrac{3\pi}{4}$.

\textbf{4.} $-\dfrac{2\pi}{3}$ : $\cos\!\left(-\dfrac{2\pi}{3}\right) = \cos\dfrac{2\pi}{3} = -\dfrac{1}{2}$ et $\sin\!\left(-\dfrac{2\pi}{3}\right) = -\sin\dfrac{2\pi}{3} = -\dfrac{\sqrt{3}}{2}$.

Coordonnées : $\left(-\dfrac{1}{2}\,;\,-\dfrac{\sqrt{3}}{2}\right)$.

%===============================================================
\section*{Exercice 2 \hfill (4 points) — C2}
%===============================================================

\baremeIndicatif{Q1 : 1,5 pt ; Q2 : 1 pt ; Q3 : 1,5 pt}

\textbf{1.} $\dfrac{5\pi}{6} = \pi - \dfrac{\pi}{6}$ :
$\cos\dfrac{5\pi}{6} = -\cos\dfrac{\pi}{6} = -\dfrac{\sqrt{3}}{2}$, \quad $\sin\dfrac{5\pi}{6} = \sin\dfrac{\pi}{6} = \dfrac{1}{2}$.

\textbf{2.} $\dfrac{7\pi}{4} = 2\pi - \dfrac{\pi}{4} = -\dfrac{\pi}{4} + 2\pi$ :
$\sin\dfrac{7\pi}{4} = -\sin\dfrac{\pi}{4} = -\dfrac{\sqrt{2}}{2}$.

\textbf{3.} $\sin^2 x = 1 - \dfrac{9}{25} = \dfrac{16}{25}$. Comme $x \in \left]\dfrac{\pi}{2};\pi\right[$, $\sin x > 0$, donc $\sin x = \dfrac{4}{5}$.

%===============================================================
\section*{Exercice 3 \hfill (4 points) — C1}
%===============================================================

\baremeIndicatif{Q1 : 2 pts ; Q2 : 1 pt ; Q3 : 1 pt}

\textbf{1.} Avec $s=r\theta$, $s=12\times\dfrac{5\pi}{6}=\boxed{10\pi\ \text{cm}}$.

\textbf{2.} $\theta=\dfrac{s}{r}=\dfrac{9}{6}=\boxed{\dfrac32\ \text{rad}}$.

\textbf{3.} La fraction de tour vaut $\dfrac{7\pi/6}{2\pi}=\boxed{\dfrac7{12}}$.

%===============================================================
\section*{Exercice 4 \hfill (4 points) — C2}
%===============================================================

\baremeIndicatif{Q1 : 1,5 pt ; Q2 : 1,5 pt ; Q3 : 1 pt}

\textbf{1.} Le point image de $-\dfrac{\pi}{3}$ a pour coordonnées $\boxed{\left(\dfrac12;-\dfrac{\sqrt3}{2}\right)}$.

\textbf{2.} Le point image de $\dfrac{3\pi}{4}$ a pour coordonnées $\boxed{\left(-\dfrac{\sqrt2}{2};\dfrac{\sqrt2}{2}\right)}$.

\textbf{3.} Oui : les coordonnées lues sur le cercle a l'angle $3\pi/4$ sont exactement celles de $M$.

%===============================================================
\section*{Exercice 5 \hfill (4 points) — C1, C2}
%===============================================================

\baremeIndicatif{Q1 : 1 pt ; Q2 : 1 pt ; Q3 : 1,5 pt ; Q4 : 0,5 pt}

\textbf{1.} La longueur parcourue vaut $25\times\dfrac{\pi}{3}=\boxed{\dfrac{25\pi}{3}\ \text{m}}$.

\textbf{2.} Le point associe a $\pi/3$ a pour coordonnées $\boxed{\left(\dfrac12;\dfrac{\sqrt3}{2}\right)}$.

\textbf{3.} La hauteur est l'altitude du centre augmentee de l'ordonnée sur la roue : $\boxed{27+\dfrac{25\sqrt3}{2}\ \text{m}}$.

\textbf{4.} $27+\dfrac{25\sqrt3}{2}\approx\boxed{48{,}7\ \text{m}}$.

\end{corrige}
